\scanpage{74}

En fait, le calcul en ``coordonnées épinglées universelles'' montre que c'est la seule cônique circonscrite (elle est donc automatiquement invariante ...)

\uline{Cas $n=3$} i.e. $\zeta=\zeta^{-2}$, $\zeta^{-1}=\zeta^{\unsure{2|1}}$, donc les fonctions pol. quasi-inv. \inserted{sous $G^+$ i.e. sous $u$} de degré $\leq2$ sont les

\[\left\{\begin{aligned}
a+bxy&&\text{val. propre }&1\\
ax+by^2&&\rule{1.6cm}{.35pt}&\zeta\\
a{\unsure{y|x}}^2+bx^2&&\rule{1.6cm}{.35pt}&\zeta^{-1}
\end{aligned}\right.\]

mais parmi ceux-ci, les qu. inv. sous $G$ sont encore $a+bxy$

\uline{Cas $n=4$} i.e. $\zeta^2=\zeta^{-2}$, les fonctions polyn. de degré $\leq2$ quasi-inv. sous $G^+$ \inserted{i.e. sous $u$} sont

\[\left\{\begin{aligned}
a+bxy&&\text{val. propre }&1\\
ax\quad(\text{resp }ay)&&\text{val. pr }&\zeta\quad(\zeta^{-1})\\
ax^2+by^2&&\text{val. propre }&\zeta^2=\zeta^{-2}
\end{aligned}\right.\]

Parmi celles-ci, les fonctions quasi-inv. sous $G$ tt entier sont les $a+bxy$, et \unsure{d'autre les|d'autre part les|distinctes} $ax^2+by^2$, \deleted{telles \unsure{telles} que}\inserted{ce qui \unsure{fait|font|forment} deux} familles à un paramètre de côniques inv. sous $G$. Dans chacune, il y a \deleted{\unsure{exactemt|exactement}} une cônique circonscrite à $S$, savoir l'habituelle $xy-1=0$, et de plus $x^2-y^2=0$

% archive sheet number: 73
